% !TeX root = ../tcc.tex
\begin{apendicesenv}
\partapendices
\chapter{Notas para revisão com os orientadores --- fora do texto do capítulo}

\textbf{TCC de Euller e Gabriel · 15/10/2026}

Versão de trabalho baseada nas leituras fornecidas. A revisão da
literatura está em andamento. O texto abaixo é uma proposta de redação;
o contexto definido é o apoio à triagem de hanseníase na rede pública de
saúde do Brasil. O GQM e a questão de pesquisa foram atualizados
conforme a definição fornecida pela dupla. As notas editoriais ao final
não integram a introdução.

\section{Estrutura adotada e relação com os
exemplos}\label{estrutura-adotada-e-relauxe7uxe3o-com-os-exemplos}

Os três TCCs fornecidos organizam suas introduções em torno de contexto,
problema, questão de pesquisa, objetivos, organização/estrutura do
trabalho e cronograma. Esta versão desenvolve contexto, problema, GQM,
questão de pesquisa, objetivos e métricas. As métricas podem ser
deslocadas para o capítulo metodológico na organização final. A
organização dos capítulos e o cronograma devem refletir o planejamento
efetivamente aprovado, sem importar atividades experimentais dos
exemplos.

Os exemplos foram utilizados como referência de organização, não como
fontes científicas sobre hanseníase ou visão computacional. Não se
reproduziu o conteúdo temático desses trabalhos.

\section{Pontos que precisam de
decisão}\label{pontos-que-precisam-de-decisuxe3o}

\begin{itemize}
\tightlist
\item
  \textbf{Recorte territorial definido:} Brasil, no contexto da rede
  pública de saúde. Isso não implica que todos os estudos devam ser
  brasileiros ou que já exista validação local. Fundamentar
  posteriormente as particularidades do cenário brasileiro com fontes
  pertinentes.
\item
  \textbf{Escopo da questão:} ``como avaliar'' pode orientar uma
  proposta de avaliação mais ampla; a revisão responde inicialmente
  ``como os modelos têm sido avaliados''. Alinhar com os professores se
  a RQ central pertence ao TCC completo e as oito perguntas orientam a
  extração da revisão.
\item
  \textbf{Tons de pele:} o escopo passa a abranger imagens de pele em
  geral. A avaliação entre tons permanece em M6, quando houver
  evidências, sem restringir a questão central a pele negra. Fitzpatrick
  não equivale a identidade racial; origem geográfica não determina tom
  de pele.
\item
  \textbf{Contribuição do TCC 1:} a versão propõe síntese da literatura
  e fundamentação de avaliação futura. Não promete execução
  experimental, validação clínica ou aplicativo concluído nesta etapa.
\item
  \textbf{Fundamentação clínica:} para aprofundar definições, sinais
  diagnósticos ou epidemiologia, consultar fontes primárias/oficiais.
  Esta versão se limita ao contexto apresentado nos artigos e evita
  estatísticas epidemiológicas atuais não verificadas.
\item
  \textbf{Lacuna provisória:} as leituras permitem apontar questões de
  avaliação e relato. Não escrever que ``não existem estudos'' ou que a
  revisão já demonstrou uma lacuna definitiva enquanto a seleção e a
  extração não terminarem.
\end{itemize}

\section{Mapa das afirmações e das
fontes}\label{mapa-das-afirmauxe7uxf5es-e-das-fontes}

{\def\LTcaptype{none} % do not increment counter
\begin{longtable}[]{@{}
  >{\raggedright\arraybackslash}p{(\linewidth - 6\tabcolsep) * \real{0.2500}}
  >{\raggedright\arraybackslash}p{(\linewidth - 6\tabcolsep) * \real{0.2500}}
  >{\raggedright\arraybackslash}p{(\linewidth - 6\tabcolsep) * \real{0.2500}}
  >{\raggedright\arraybackslash}p{(\linewidth - 6\tabcolsep) * \real{0.2500}}@{}}
\toprule\noalign{}
\begin{minipage}[b]{\linewidth}\raggedright
Parte da redação
\end{minipage} & \begin{minipage}[b]{\linewidth}\raggedright
Base consultada
\end{minipage} & \begin{minipage}[b]{\linewidth}\raggedright
Onde conferir
\end{minipage} & \begin{minipage}[b]{\linewidth}\raggedright
Limite da afirmação
\end{minipage} \\
\midrule\noalign{}
\endhead
\bottomrule\noalign{}
\endlastfoot
Dificuldades de reconhecimento e panorama de IA & \citeonline{fernandes2024} & PDF pp.~1--2 & Contextualização da revisão, sem números
atuais de prevalência \\
21 estudos e predominância de classificação multiclasse & \citeonline{fernandes2024} & Resumo, p.~1 & Não significa 21 estudos elegíveis para o
protocolo de vocês \\
MobileNetV2 e protótipo Android & \citeonline{jaikishore2021} &
Resumo/Introdução, p.~1 & Protótipo não equivale a validação clínica;
não afirmar estadiamento validado da hanseníase \\
CNNs e população de pele escura & \citeonline{yotsu2023} & pp.~3--5 & Dados
africanos não demonstram validação na rede pública brasileira \\
Resultados específicos de hanseníase & \citeonline{yotsu2023} & pp.~6--7,
Task 2/Figura 1 & 12\%/13\% são leitura da matriz por classe; não
misturar com Task 1 \\
Segmentação de feridas & \citeonline{sanchez2024} & pp.~1 e 4--5 do PDF &
Não é diagnóstico da hanseníase \\
YOLO e explicabilidade & \citeonline{eita2025} & pp.~4--9 & Mapas de
saliência não comprovam validade clínica \\
Diversidade limitada e validação prospectiva futura & \citeonline{eita2025}
& p.~16, seções 6.4 e 7 & Não demonstra diferença quantitativa por tom
de pele \\
\end{longtable}
}

A classificação de \citeonline{ahikpa2025} como revisão ou estudo experimental requer confirmação no texto completo. Pode apoiar o referencial e os trabalhos
relacionados; nesta introdução, priorizaram-se a revisão sistemática de
Fernandes e os estudos primários fornecidos. Isso evita atribuir aos
revisores resultados que pertencem aos estudos citados.

\section{Situação do repositório
consultado}\label{situauxe7uxe3o-do-reposituxf3rio-consultado}

Consulta à branch \texttt{master} de \texttt{HansenSightAI/tcc}, em
14/10/2026. O arquivo \texttt{Monografia/editaveis/introducao.tex} ainda
contém uma introdução sobre análise de shaders em dispositivos móveis. O
arquivo \texttt{Monografia/bibliografia.bib} contém referências do mesmo
modelo de computação gráfica. Esta redação é um rascunho independente
para revisão; não foi aplicada ao repositório.

Caminhos consultados:

\begin{itemize}
\tightlist
\item
  https://github.com/HansenSightAI/tcc/blob/master/Monografia/editaveis/introducao.tex
\item
  https://github.com/HansenSightAI/tcc/blob/master/Monografia/bibliografia.bib
\end{itemize}

Na integração, substituir o texto do modelo pela versão aprovada e
acrescentar as referências abaixo. Antes de remover referências antigas,
verificar se ainda são citadas em outros capítulos do modelo.


\chapter{Solicitação enviada ao Consensus}
\label{ap:prompt-consensus}

O texto a seguir reproduz a solicitação enviada em inglês à ferramenta
para orientar a busca. As quebras de linha foram ajustadas à apresentação;
os termos e os asteriscos foram preservados.

\lstinputlisting[basicstyle=\small\ttfamily,breaklines=true,breakatwhitespace=true,columns=fullflexible]{editaveis/prompt-consensus.txt}

\end{apendicesenv}
